% Section commune : détection d'objets sur FPGA
\section{Détection d'objets sur FPGA}
\label{sec:detection}

\subsection{Trajectoire du corpus}

Les premiers travaux sur la détection partent d'un constat de
\citet[§1]{2019-ding-req-yolo-resource-aware-efficient} : la littérature FPGA accélère
surtout la classification. En 2018, aucun toolflow ne visait la détection
\citep[§2.6, §3.1]{2018-venieris-toolflows-mapping-convolutional-neural-networks}. Le
corpus montre ensuite quatre approches.

\paragraph{Compression puis design sur mesure.} REQ-YOLO compresse Tiny YOLO (YOLOv2-tiny)
de deux façons :
\begin{itemize}
  \item des matrices bloc-circulantes, qui transforment la convolution en FFT, produit
    terme à terme et FFT inverse \src{§3} ;
  \item une quantification hétérogène par ADMM qui choisit, couche par couche, entre des
    multiplications sur DSP et des décalages sur LUT \src{§4}.
\end{itemize}
Tous les poids tiennent alors en BRAM \src{§5}. Sur un Virtex-7 690T en carte PCIe, le
modèle atteint 314,2\,FPS pour 21\,W \src{Table~2}. Trois réserves :
\begin{itemize}
  \item la latence de 3,18\,ms est une somme des latences par couche, sans transfert PCIe
    ni NMS \src{Table~1} ;
  \item la mAP absolue n'est pas dans le texte : seules des pertes relatives le sont
    \src{§6.2} ;
  \item la plateforme n'a rien d'embarqué.
\end{itemize}
Face à une Jetson TX2 exécutant le même modèle (28,4\,FPS, 10,8\,W, \src{Table~2}), le
gain est net, mais TensorRT n'est pas mentionné : la TX2 est vraisemblablement
sous-exploitée \deduit.

\paragraph{Protocole commun, mesuré par un tiers.} Le concours DAC-SDC 2018
\citep{2021-xu-dac-sdc-low-power-object} impose deux plateformes (PYNQ-Z1 et Jetson TX2),
un jeu de test caché d'images de drone, un débit minimal et une énergie mesurée au
wattmètre par les organisateurs \src{§2.3}. C'est le seul protocole commun FPGA/GPU du
corpus. Les équipes FPGA gagnantes utilisent des détecteurs maison en HLS ou RTL sur
XC7Z020, avec les cartes d'activation en DDR \src{§4.2} :
\begin{itemize}
  \item TGIIF, un SSD élagué en 8~bits : IoU de 0,62 à 11,96\,FPS \src{§5.2.1} ;
  \item SystemsETHZ, un SqueezeNet binaire suivi d'une tête YOLO : IoU de 0,49 à
    25,97\,FPS \src{§5.2.1}.
\end{itemize}
Limites : la métrique est l'IoU d'un seul objet par image, pas une mAP ; les puissances
par équipe ne figurent qu'en figures ; le score composite utilise des bases de logarithme
différentes pour le FPGA et le GPU \src{§2.3}. Les auteurs estiment que le FPGA consomme
1/3 à 1/2 de l'énergie du GPU pour une précision proche, sans chiffre dans le texte
\src{§5.3}.

\paragraph{Design HLS sur SoC d'entrée de gamme.}
\citet{2023-zhai-fpga-vehicle-detection-tracking-accelerator} élaguent YOLOv3-tiny par
canaux (calcul divisé par 7,4), le quantifient en virgule fixe 16~bits dynamique et
l'implantent en HLS sur un Zynq XC7Z035. Leur accélérateur combine Winograd pour les
convolutions 3$\times$3 et im2col + GEMM pour les 1$\times$1 \src{§3.1, §3.5}. Le NMS et
le suivi Deepsort tournent sur l'ARM \src{§3.3}. Avec un module, ils obtiennent
91,65\,FPS pour 12,51\,W ; avec deux modules, 168,72\,FPS pour 15,18\,W
\src{Tables~8--9}. La méthode de mesure de puissance n'est pas décrite. Deux points
critiques :
\begin{itemize}
  \item l'AP@0.5 est identique en FP32 et en INT16 à trois décimales, ce qui laisse penser
    qu'elle n'a pas été remesurée sur la cible \deduit ;
  \item la ligne Jetson Nano de leur comparaison vient d'un autre article et porte sur un
    modèle non élagué \src{Table~7}. Le gain affiché vient donc surtout de l'élagage.
\end{itemize}

\paragraph{Toolflow streaming et flot industriel.} SATAY
(section~\ref{sec:toolflows}) couvre YOLOv3-tiny à YOLOv8s, mais ses meilleurs chiffres
viennent de cartes Virtex de 23 à 67\,W \src{Table~III}. Sur la seule carte SoC testée,
la ZCU104, YOLOv5n en 640$\times$640 obtient 26,7 de mAP50 COCO pour 21,41\,ms et
14,5\,W, contre 27,7 de mAP50, 32,28\,ms et 8,58\,W pour une Jetson TX2
\src{Table~IV}$^{m}$. L'énergie par inférence est du même ordre (317 contre 277\,mJ,
\src{Table~IV}). À l'autre extrémité,
\citet{arxiv-2507.18174-real-time-object-detection-classification} déploient YOLOv5 INT8
par Vitis~AI sur KV260 : 9\,FPS et 3,5\,W, méthode non décrite \src{Table~I}. Leur
précision (98,6\,\%) est une \enquote{classification accuracy} sur 72 images choisies, qui
ne se compare à aucune mAP. Leur propre tableau contredit leur conclusion : 2,6\,FPS/W,
sous quatre des cinq travaux cités \src{§IV-B, Table~I}.

\subsection{Tableau des détecteurs FPGA du corpus}

\begin{table}[!htbp]
\centering\footnotesize
\caption{Détecteurs sur FPGA dans \texttt{benchmarks.csv}. Précisions non comparables
d'une ligne à l'autre (jeux, métriques et résolutions différents). $^{m}$ : confiance
\emph{medium}. Portée de puissance : u = \emph{unspecified}, b = \emph{board}.}
\label{tab:detection-fpga}
\begin{tabularx}{\textwidth}{@{}P{2.2cm}LP{1.55cm}P{2.1cm}rrrl@{}}
\toprule
Travail & Modèle, précision & Puce & Métrique & Préc. & FPS & ms & W (portée) \\
\midrule
REQ-YOLO & Tiny YOLO bloc-circulant, bits \nr & XC7VX690T & mAP VOC & \nr & 314,2 & 3,18 & 21 (u) \\
FINN-R & Tincy YOLO W1A3 (MO) & XC7Z020 & mAP VOC & 50,1 & \nr & \nr & 2,5 (b) \\
FINN-R & Tincy YOLO W1A3 (MO) & XCZU3EG & mAP VOC & 50,1 & \nr & \nr & 9,7 (b) \\
DAC-SDC & TGIIF : SSD élagué INT8 & XC7Z020 & IoU, 1 objet & 62,4 & 11,96 & \nr & \nr$^{m}$ \\
DAC-SDC & SystemsETHZ : SqueezeNet binaire + YOLO & XC7Z020 & IoU, 1 objet & 49,2 & 25,97 & \nr & \nr$^{m}$ \\
Zhai & YOLOv3-tiny élagué INT16, 1 module & XC7Z035 & AP@0.5, 1 classe & 59,9 & 91,65 & \nr & 12,51 (u) \\
Zhai & idem, 2 modules & XC7Z035 & --- & \nr & 168,72 & \nr & 15,18 (u) \\
SATAY & YOLOv5n 640, PTQ & ZCU104 & mAP50 COCO & 26,7 & \nr & 21,41 & 14,5 (u)$^{m}$ \\
SATAY & YOLOv8s 640, W8A16 & VCU118 & mAP$^{*}$ COCO & 61 & \nr & 24,5 & 57,4 (u)$^{m}$ \\
Magalhães & RetinaNet R50 INT8, 2 DPU & ZCU104 & mAP VineSet & \nr & 25 & \nr & \nr$^{m}$ \\
Al Amin & YOLOv5 INT8 (DPU) & KV260 & acc., 72 images & 98,6 & 9 & \nr & 3,5 (u) \\
\bottomrule
\end{tabularx}
\end{table}

Sources : REQ-YOLO \citep[Tables~1--2]{2019-ding-req-yolo-resource-aware-efficient} ; FINN-R
\citep[Table~5]{arxiv-1809.04570-finn-r-end-end-deep} ; DAC-SDC
\citep[§5.2.1]{2021-xu-dac-sdc-low-power-object} ; Zhai
\citep[Tables~8--9]{2023-zhai-fpga-vehicle-detection-tracking-accelerator} ; SATAY
\citep[Tables~III--IV]{2023-montgomerie-corcoran-satay-streaming-architecture-toolflow-accelerating} ;
Magalhães \citep[§3]{arxiv-2211.11647-benchmarking-edge-computing-devices-grape} ; Al Amin
\citep[Table~I]{arxiv-2507.18174-real-time-object-detection-classification}. La latence de REQ-YOLO est une somme des
latences par couche. Le mAP$^{*}$ de SATAY porte un astérisque dont la définition n'a pas
été extraite.

\paragraph{Lecture du tableau.} D'après un recompte de \texttt{benchmarks.csv}, sur 29
lignes de détection sur FPGA, 14 ont une mAP COCO, et toutes viennent de SATAY. Aucune ligne
ne rapporte à la fois une mAP standard, un FPS de bout en bout et une puissance de carte
mesurée. Les surveys de 2025 confirment ce manque d'homogénéité :
\begin{itemize}
  \item \citet{2025-hozhabr-survey-real-time-object-detection} proposent de normaliser par
    le \emph{pixel throughput} (FPS $\times$ pixels par image, \src{§II-C2}). Cela corrige
    la résolution, mais ni la précision ni la portée de puissance ;
  \item \citet{arxiv-2509.04153-real-time-fpga-cnns-detection}\pp{} alignent des designs
    mesurés dans des conditions hétérogènes, avec des incohérences entre texte et tableaux
    (voir la fiche).
\end{itemize}
Hozhabr et Giorgi affirment que \enquote{FPGAs have emerged as the optimal hardware
platform, offering both speed and power efficiency} \src{§IV}, sans comparaison mesurée
avec une Jetson récente ni un MCU à NPU.

\subsection{Post-traitement et interaction PS--PL}

Le NMS, le suivi et le redimensionnement tournent presque toujours sur le processeur du SoC
et sont rarement comptés :
\begin{itemize}
  \item Zhai et al. : NMS et Deepsort sur l'ARM, inclus ou non dans le FPS, on ne sait pas
    \src{§3.3} ;
  \item REQ-YOLO : latence sans NMS \src{Table~1} ;
  \item Magalhães et al. : cœur du réseau seul, pré- et post-traitement exclus
    \src{§2.5} ;
  \item DAC-SDC : redimensionnement sur l'ARM \src{§4.2}.
\end{itemize}
Le survey de Sali et al. cite un cas où le débit passe de 195\,FPS pour la détection seule
à 24\,FPS de bout en bout, à cause d'un post-traitement en Python \src{§VI-C} (chiffre d'un
travail cité, non vérifié). Hozhabr et Giorgi citent deux contre-exemples, avec un NMS
pipeliné en matériel ou un système entièrement sur FPGA \src{§V-A}. Tous deux sont sur de
grosses puces (Stratix~10, Kintex-7) et restent à lire à la source.
